\problema{Lowest Common Ancestor of a Binary Tree II}{1644}{src/02-arboles/1644-lca-ii.cpp}{M}

Es la variante del problema 236 donde \textbf{p y/o q pueden NO existir} en
el árbol. Si alguno de los dos falta, debemos devolver \texttt{nullptr} en
vez de un ancestro cualquiera.

\paragraph{Intuición.} En 236 podíamos cortar la búsqueda apenas
encontrábamos p o q, porque el enunciado garantizaba que ambos vivían en el
árbol: encontrar uno era suficiente para confiar en el resultado. Aquí esa
garantía desaparece. Si cortáramos temprano y el otro nodo resultara no
existir, devolveríamos un LCA falso. La solución: recorrer el árbol
\emph{completo}, sin atajos, y llevar un contador de cuántos de los dos
nodos objetivo se encontraron de verdad. Sólo si el contador llega a 2
confiamos en el candidato a LCA; si es 0 o 1, alguno falta y respondemos
\texttt{nullptr}.

\begin{center}
\ttfamily
\begin{tabular}{ccc}
     & 3 &     \\
  5  &   &  1  \\
6\ 2 &   & 0\ 8\\
     &7\ 4&    \\
\end{tabular}
\end{center}

\noindent Con este árbol: \texttt{LCA(5, 1) = 3}, \texttt{LCA(5, 4) = 5},
pero \texttt{LCA(5, 99) = nullptr} porque 99 no está en el árbol, y
\texttt{LCA(99, 88) = nullptr} porque ninguno de los dos existe.

\paragraph{Solución.} Usamos el mismo DFS post-order de 236 (\texttt{dfs}
devuelve, para cada subárbol, el candidato a LCA), pero SIEMPRE evaluamos
ambos hijos antes de revisar el nodo actual (sin retorno anticipado), y cada
vez que el nodo actual es p o q incrementamos un contador pasado por
referencia. Al final, el resultado sólo es válido si el contador es
exactamente 2:

\lstinputlisting{src/02-arboles/1644-lca-ii.cpp}

\complejidad{$O(n)$}{$O(h)$}

\paragraph{Notas de entrevista (Amazon).} Deja claro en voz alta la
diferencia con 236: aquí no hay garantía de existencia, así que el corte
temprano (\texttt{if root == p or root == q: return root}) es incorrecto
porque puede devolver un LCA falso cuando el otro nodo no está. La técnica
del contador por referencia es la forma limpia de confirmar existencia sin
un segundo recorrido; también es válida una alternativa de dos pasadas (una
para verificar que p y q existen, otra para calcular el LCA como en 236),
pero es menos elegante y hace doble trabajo. Casos borde: p o q ausentes,
ambos ausentes, p es ancestro de q (o viceversa), y árbol de un solo nodo.
